\Qhead{1.4}{Аналитические и качественные методы теории нелинейных колебаний}{Аналитические и качественные методы теории нелинейных колебаний. Метод гармонического баланса. Укороченные уравнения. Усреднение в системах, содержащих быстрые и медленные движения. Анализ возможных движений и бифуркаций в фазовом пространстве.}
\Subfile{1.4.1}{sections/01_oscillations/q04/s01}
\Subfile{1.4.2}{sections/01_oscillations/q04/s02}
\Subfile{1.4.3}{sections/01_oscillations/q04/s03}
\Subfile{1.4.4}{sections/01_oscillations/q04/s04}

\recap{Гармонический баланс ищет периодический режим; МММА --- медленные $A(t),\phi(t)$ и устойчивость. Усреднение объясняет действие медленного изменения и быстрой вибрации. Фазовый портрет показывает движения и их бифуркации.}
\board{Первая гармоника $x^3$; две проекции МММА; потенциал Капицы; узел, фокус, седло, центр и цикл.}
\oral{Почему при усреднении быстро осциллирующей силы недостаточно заменить её нулём?}
\bait{Можно оставить адиабатический инвариант: при медленной перестройке сохраняется $E/\omega$, потому что привод меняет энергию.}
\sources{Пятаков~\cite{pyatakov}, PDF-с.~10--23, 32--42, 50--53: фазовые методы, баланс, разделение времён и МММА. Проекция на косинус восстановлена самостоятельно вместо опечатки PDF-с.~53. Подача методов переработана по вопросу~4 коллеги; нормальные формы бифуркаций --- самостоятельное математическое пояснение.}
